\documentclass[11pt, a4paper]{article}

\usepackage[utf8]{inputenc}
\usepackage[T1]{fontenc}
\usepackage{geometry}
\geometry{margin=1in}
\usepackage{titlesec}
\usepackage{amsmath, amssymb}
\usepackage{booktabs}
\usepackage{tabularx}
\usepackage{hyperref}
\usepackage{enumitem}

\hypersetup{
    colorlinks=true,
    linkcolor=blue!70!black,
    urlcolor=blue!70!black,
    pdftitle={Genus 2 Curve Arithmetic System - Technical Overview}
}

\title{\textbf{Genus 2 Curve Arithmetic System} \\ \large Technical Overview}
\author{}
\date{}

\begin{document}

\maketitle

\begin{abstract}
Rational point search $\cdot$ Jacobian rank computation $\cdot$ Finite-field DLP.
\end{abstract}

\section{Executive Summary}

This system is a research-grade SageMath/Python toolkit for studying the arithmetic of genus 2 (and higher-genus) hyperelliptic curves over the rationals. Its central innovation is an iterated fibration tower construction that reduces difficult questions about genus $\geq 2$ curves to tractable problems on families of elliptic curves. Building on this foundation, the system provides three tightly integrated capabilities:

\begin{itemize}[label={\textbullet}]
    \item Rational point enumeration — finding rational points on curves of arbitrary genus (genus 2 fully supported; genus 3 tested; higher genera possible) starting from a single seed point, using CRT, LLL lattice reduction, and Hensel lifting.
    \item Jacobian rank determination (\texttt{MUMFORD\_SEARCH} mode) — computing a provably independent basis for $J(\mathbb{Q})$, the Mordell--Weil group of the Jacobian, typically in well under two minutes regardless of conductor or rank.
    \item Finite-field DLP / index calculus — mounting a genus 2 hyperelliptic-curve discrete-logarithm attack via Mumford divisor smoothness relations and sparse linear algebra.
\end{itemize}

The system has been validated against the LMFDB on dozens of genus 2 curves spanning a wide range of conductors (up to $\sim 10^7$) and Mordell--Weil ranks (0 through 4+).

\section{Mathematical Background}

\subsection{Genus 2 Curves and Their Jacobians}

A genus 2 curve $C$ over $\mathbb{Q}$ is a smooth projective curve of the form $y^2 = f(x)$ where $f$ is a squarefree polynomial of degree 5 or 6. Unlike elliptic curves, $C(\mathbb{Q})$ need not itself form a group; instead, its degree-0 Picard group --- the Jacobian $J(C)$ --- is a 2-dimensional abelian variety that generalizes the group law. By the Mordell--Weil theorem, $J(\mathbb{Q})$ is a finitely generated abelian group. Its free rank $r = \text{rank}\, J(\mathbb{Q})$ (often just called the rank of $C$) is a central arithmetic invariant. Computing $r$ rigorously is much harder for genus 2 than for elliptic curves, where standard descent methods apply directly.

Elements of $J(\mathbb{Q})$ are represented as Mumford divisors: pairs $(u(x), v(x))$ where $u$ is monic of degree 2, $v$ has degree $< 2$, and $u \mid v^2 - f(x)$ in $\mathbb{Q}[x]$. The Cantor algorithm handles group-law computations in this representation.

\subsection{Elliptic Fibrations}

An elliptic fibration is a map $\pi: S \to \mathbb{P}^1$ from a surface $S$ to the projective line whose generic fiber is a smooth elliptic curve $E/\mathbb{Q}(m)$. The system constructs such a fibration starting from the genus 2 curve $C$ and a rational point on $C$. The parameter $m$ lives in the intersection locus: it records where the tangent (or secant) line through a known rational point meets the shifted curve. A rational $m$-value corresponds to a rational point on $C$, which is precisely the information the search harvests via CRT.

\subsection{Shioda--Tate and Picard Numbers}

For an elliptic surface $S/\mathbb{Q}$ with generic fiber $E/\mathbb{Q}(m)$, the Shioda--Tate formula relates the Picard number $\rho(S)$, the rank of $E(\mathbb{Q}(m))$, and the fiber contribution:
$$\rho(S) = 2 + \text{rank}\, E(\mathbb{Q}(m)) + \sum_v (m_v - 1)$$
where the sum is over singular fibers of Kodaira type, and $m_v$ is the number of irreducible components. The system computes both a lower bound on $\rho$ (from characteristic-0 geometry) and an upper bound (from reduction modulo good primes via the Lefschetz trace formula), often pinning $\rho$ exactly and thereby determining the Mordell--Weil rank of the fibration.

\section{Core Algorithm: The Fibration Tower}

\subsection{Tower Construction}

Given a genus 2 curve $C: y^2 = f(x)$ and a seed rational point $P = (x_0, y_0) \in C(\mathbb{Q})$, the system constructs an iterated fibration tower:

\begin{itemize}[label={\textbullet}]
    \item Degree 5 polynomials (genus 2, one infinite place): one tower step reduces to a degree-4 (genus 1) quartic fibration.
    \item Degree 6 polynomials (genus 2, two infinite places): two tower steps are required, first reducing to degree 5, then to degree 4.
\end{itemize}

At each step the algorithm:
\begin{enumerate}
    \item Selects the best-scoring auxiliary polynomial $Q(x)$ of appropriate degree using a complexity measure (height, degree, collision resistance of the discriminant locus).
    \item Computes the intersection locus as a rational function $r(m)$ in the fibration parameter $m$.
    \item Passes to the next step using the residual degree-4 quartic as the new fiber.
\end{enumerate}

Scores are computed across 10 candidate geometries and the lowest-complexity choice is selected, trading off coefficient height, degree, and bad-prime content. Non-minimal models are sometimes preferred despite seemingly higher complexity, because they are less sensitive to height bounds.

\subsection{Weierstrass Model and Fiber Analysis}

Once the quartic fiber is established, the system computes the Weierstrass model $E: y^2 = x^3 + a_4(m)x + a_6(m)$ over the function field $\mathbb{Q}(m)$. This involves:

\begin{itemize}[label={\textbullet}]
    \item Minimal Weierstrass reduction: Tate's algorithm applied fiber by fiber over algebraic places of $m$.
    \item Kodaira fiber classification: each irreducible factor of the discriminant $\Delta(m)$ is classified ($I_n$, $II$, $III$, $IV$, $I^*_n$, $\dots$) by examining the valuations of $c_4$, $c_6$, $\Delta$.
    \item Euler characteristic validation: the sum of fiber contributions must equal 12 for a standard rational elliptic surface.
    \item Discriminant degree check: degree 12 in $m$ indicates a fully minimal model; deviations trigger diagnostics.
\end{itemize}

\subsection{Mordell--Weil Sections}

From the known rational point $P$, the system generates sections --- rational functions $S_i(m) = (x_i(m), y_i(m))$ satisfying the Weierstrass equation for all $m$. Sections are the Mordell--Weil group elements of the generic fiber. Multiple independent sections arise from multiples $[n]P$ of the base section and from additional rational points found during the search. Independence is verified via the N\'eron--Tate height pairing matrix, which must be positive definite for a true basis.

\section{Rational Point Search over $\mathbb{Q}$}

\subsection{The Search Problem}

The key observation is that a rational point $(x_0, y_0) \in C(\mathbb{Q})$ corresponds to a rational value of $m$ satisfying $x_0 = r(m)$ where $r(m)$ is the intersection-locus rational function. The search therefore reduces to finding rational roots of the numerator of $S_i(m) - r(m)$ for sections $S_i$.

\subsection{Modular Arithmetic Pipeline}

The system searches for rational $m$ by collecting residues modulo a pool of primes and applying CRT reconstruction:

\begin{enumerate}
    \item For each prime $p$ in the pool, solve $S_i(m) \equiv r(m) \pmod{p}$ to find candidate residues $m \equiv c \pmod{p}$.
    \item Apply Hensel's lemma to lift simple roots to higher precision.
    \item Filter by torsion order (small-order residues likely correspond to torsion sections, not rational points).
    \item Combine residues across multiple prime subsets via CRT, forming candidate rational numbers by rational reconstruction.
    \item Verify each candidate by substituting back into the curve equation.
\end{enumerate}

\subsection{LLL Lattice Reduction}

Raw search vectors (integer combinations of sections) are preprocessed by LLL and BKZ lattice reduction to find short vectors. Short vectors correspond to low-multiplier combinations $[n]P$ that are most likely to produce small-height rational $m$ values, dramatically reducing the search space. Configurable parameters include the height bound, prime pool size, LLL delta, BKZ block size, and extra-prime filtering density. The system auto-tunes several of these from the Galois group of the discriminant and the canonical height of the seed point.

\subsection{Multi-Fibration Strategy}

Different fibrations --- built from different auxiliary polynomials $Q(x)$ --- expose different rational points because each section sees the curve's geometry from a different angle. Not every fibration can detect every rational point. The system therefore runs multiple fibrations sequentially, adding newly discovered points as seeds for subsequent fibration towers. A consensus filter can optionally require agreement across several fibrations before accepting a candidate.

\section{Jacobian Rank Finding: MUMFORD\_SEARCH Mode}

This mode directly computes the Mordell--Weil rank of $J(C)$ by constructing an explicit independent basis for $J(\mathbb{Q})$. It runs in well under 2 minutes for genus 2 curves of arbitrary conductor and rank.

\subsection{What It Computes}

Whereas the rational point search looks for rational $m$-values (i.e., rational points on $C$), \texttt{MUMFORD\_SEARCH} mode looks for rational Mumford divisors: elements $D \in J(\mathbb{Q})$ of the form $[u(x), v(x)]$ with $u, v \in \mathbb{Q}[x]$, $\deg u = 2$, $\deg v < 2$. The size of the largest independent such set is the Mordell--Weil rank.

\subsection{The Mumford Residue System}

Each Mumford divisor is specified by four rational numbers $(s, p, v_1, v_0)$, where $u(x) = x^2 - sx + p$ and $v(x) = v_1 x + v_0$. The Mumford conditions give a system of polynomial constraints. The fibration provides a modular parameterization: for each prime $\ell$ in the pool, the system solves the 5-equation Mumford system over $\mathbb{F}_\ell$ to harvest candidate residues of $(s, p, v_1, v_0)$.

\subsection{CRT Reconstruction of Divisors}

Residues of the four Mumford coordinates are collected across the prime pool and combined via CRT. The reconstruction phase:

\begin{itemize}[label={\textbullet}]
    \item Groups vectors by their LLL-reduced section multiple, focusing reconstruction on vectors most likely to lift to rational divisors.
    \item Applies algebraic pre-filtering using the consistency equations to prune incompatible residue combinations before CRT.
    \item Performs height-filtered rational reconstruction, accepting only candidates whose Weil height is within the configured bound.
    \item Verifies each reconstructed $(s, p, v_1, v_0)$ against the Mumford conditions over $\mathbb{Q}$.
\end{itemize}

\subsection{Independent Basis Construction}

After reconstruction, the system selects an independent basis from the candidate divisors using a multi-prime certification protocol:

\begin{itemize}[label={\textbullet}]
    \item Torsion filter: divisors with finite order in $J(\mathbb{Q})$ are identified by checking $[N]D = 0$ for small $N$ using the group law, and discarded.
    \item Mod-$p$ scoring: each divisor is evaluated modulo several good primes; those with consistently high ``rank score'' across primes are prioritized.
    \item Incremental certification: divisors are added to the basis one at a time. A new divisor is accepted only if it is independent of the current basis modulo at least 3 distinct primes. Independence is checked via Smith normal form on the relation matrix.
    \item Height pairing confirmation: once the basis is certified by mod-$p$ methods, the N\'eron--Tate Gram matrix $H = (\langle D_i, D_j \rangle)$ is computed using Arakelov height pairings and confirmed to be positive definite.
\end{itemize}

The rank is simply the final basis size. Because the certification uses provably correct linear algebra over $\mathbb{Z}/n\mathbb{Z}$ (Smith normal form) at multiple primes, the result is certified rather than heuristic.

\subsection{Typical Performance}

For genus 2 curves from the LMFDB with conductors up to $\sim 10^7$ and Mordell--Weil ranks 0 to 4, the complete rank-finding computation --- fibration construction, residue harvesting, CRT reconstruction, torsion filtering, and basis certification --- typically completes in under 2 minutes on a modern multicore workstation. The bottleneck is the CRT reconstruction phase, which is parallelized across up to 20 workers.

As an example, a rank-4 curve ($y^2 = x^6 + 8x^5 + 10x^4 - 10x^3 - 11x^2 + 2x + 1$) with 11 known rational points: the system found a certified basis of 4 independent Mumford divisors and all 7 non-trivial rational $x$-coordinates in 33.5 seconds total.

\section{Finite-Field Index Calculus (DLP Mode)}

\subsection{Problem Setup}

When \texttt{FINITE\_FIELD = p} is set, the system attacks the genus 2 Hyperelliptic Curve Cryptosystem (HECC) discrete logarithm problem: given $G, Q \in J(\mathbb{F}_p)$ with $Q = [d]G$, find $d$. The system automatically generates a cryptosystem instance: it constructs a prime-order subgroup of $J(\mathbb{F}_p)$, picks a random base divisor $G$, computes $Q = [d]G$ for a chosen secret key $d$, and records four preferred $x$-coordinates (the supports of $G$ and $Q$) to aid smoothness detection.

\subsection{Index Calculus Pipeline}

The attack proceeds in three stages:

\begin{enumerate}
    \item Factor base construction: smooth Mumford divisors are collected using the fibration-parameterized Mumford residue system over $\mathbb{F}_p$. ``Smooth'' means $u(x)$ splits completely over $\mathbb{F}_p$ into linear factors, giving degree-1 components that form the factor base.
    \item Relation generation: the random walk harvests divisors expressible as sums of factor-base elements. A structured walk (seeded by $x$-residues from the fibration) and a collision walk (Baby-step Giant-step) cooperate to reach the target relation count.
    \item Linear algebra: the relation matrix (sparse, over $\mathbb{Z}/\ell\mathbb{Z}$ where $\ell$ is the group order) is solved using Block Wiedemann's algorithm. The solution gives $d$ directly.
\end{enumerate}

\subsection{Matrix Rank Diagnostics}

Before committing to linear algebra, the system verifies that the relation matrix is full rank. A full-rank matrix guarantees a unique solution; rank deficiency indicates insufficient or linearly dependent relations and triggers further relation collection.

\section{Supporting Analyses}

\subsection{Picard Number ($\rho$) via Shioda--Tate}

Using the Shioda--Tate formula, the system computes both a characteristic-0 lower bound on $\rho(S)$ (from the known sections and fiber types) and an upper bound from the rank of the N\'eron--Severi group modulo a good prime $\ell$ (via the Lefschetz trace on the $\ell$-adic cohomology). When both bounds agree, $\rho$ is determined exactly, which in turn pins the Mordell--Weil rank of the generic fiber.

\subsection{2-Selmer Rank Bounds}

A 2-descent computes an upper bound on $\text{rank}\, J(\mathbb{Q})$ via the Selmer group. Local conditions are evaluated at each bad prime (including $\infty$) and the global Selmer rank is bounded by the sum of local dimensions minus the reciprocity correction. This provides an upper bound complementing the lower bound from the Mumford basis.

\subsection{Torsion Analysis}

The system identifies the torsion subgroup of $J(\mathbb{Q})$ by computing the GCD of specializations of the group order at good primes and by checking small multiples $[N]D = 0$. Torsion elements are filtered out before basis construction to ensure the Mumford basis is free.

\subsection{Automorphism Analysis}

The N\'eron--Severi lattice automorphisms (scaling and translation symmetries of the fibration) are computed and reported. Non-trivial automorphisms can be exploited to reduce the search space or to relate different fibrations.

\subsection{CM Fiber Detection}

Special fibers with complex multiplication ($j$-invariant $0, 1728, -19728, -88736, \dots$) are detected by solving $a_4(m) = 0$ or specialized $j$-invariant equations. CM fibers often yield high-rank specializations and are tried as candidate rational points.

\subsection{Archimedean Height Pairings}

The N\'eron--Tate height pairing $\langle D_i, D_j \rangle$ between Mumford divisors is computed using a combination of Arakelov intersection theory (for the finite places) and numerical integration over the Jacobian's period lattice (for the archimedean place). These pairings form the Gram matrix $H$, whose determinant and positive-definiteness certify independence of the basis.

\section{Feature Summary}

\begin{table}[h]
    \centering
    \renewcommand{\arraystretch}{1.2}
    \begin{tabularx}{\textwidth}{l X}
        \toprule
        \textbf{Feature} & \textbf{Description} \\
        \midrule
        Fibration Tower & Iterated 1--2 step tower reducing genus 2 curve to elliptic fibration over $\mathbb{Q}(m)$; best-of-10 geometry selection. \\
        Weierstrass \& Kodaira & Full Weierstrass minimization, Tate algorithm, Kodaira fiber classification, Euler characteristic. \\
        Rational Point Search & LLL + CRT + Hensel + rational reconstruction pipeline; auto-configured height bounds and prime pool. \\
        Rank Finding (MUMFORD) & Provably independent Mumford basis for $J(\mathbb{Q})$; mod-$p$ certification + Smith normal form; under 2 min typical. \\
        Torsion Filtering & Identifies and removes torsion Mumford divisors before basis construction. \\
        Picard / Shioda--Tate & Exact $\rho(S)$ via char-0 lower bound + reduction upper bound; determines MW rank of generic fiber. \\
        2-Selmer Bounds & Local 2-descent at all bad primes; upper bound on $\text{rank}\, J(\mathbb{Q})$. \\
        CM Fiber Detection & Special $j$-invariant fiber search for high-rank specializations. \\
        Automorphisms & NS lattice translation and scaling automorphisms; symmetry exploitation. \\
        Finite-Field DLP & HECC index calculus: smooth divisor factory, structured + collision walk, Block Wiedemann solve. \\
        Height Pairings & Arakelov + archimedean N\'eron--Tate pairings; positive-definite Gram matrix enforced. \\
        Multi-Fibration & Sequential fibration construction from growing point list; consensus filter across fibrations. \\
        Diagnostics & Completeness posterior, prime MI analysis, saturation ratio, sufficiency proofs. \\
        Parallel Execution & Residue computation and CRT reconstruction parallelized across up to 20 workers. \\
        \bottomrule
    \end{tabularx}
    \caption{System Capabilities and Features}
\end{table}

\section{Usage and Configuration}

\subsection{Primary Entry Point}

The system is invoked via the main entry point. All configuration lives in the configuration file.

\subsection{Key Configuration Flags}

\begin{table}[h]
    \centering
    \renewcommand{\arraystretch}{1.2}
    \begin{tabularx}{\textwidth}{l X}
        \toprule
        \textbf{Flag} & \textbf{Description} \\
        \midrule
        COEFFS\_GENUS2 & Coefficient list $[a_n, \dots, a_0]$ defining $y^2 = f(x)$. Degree 5 or 6 for genus 2; degree 7--8 for genus 3; higher degrees accepted (experimental). \\
        DATA\_PTS\_GENUS2 & List of seed $x$-coordinates (rational) on the curve. \\
        TERMINATE\_WHEN\_6 & Stop searching after finding this many distinct rational $x$-coordinates. \\
        MUMFORD\_SEARCH & \texttt{True}: find Mumford basis for $J(\mathbb{Q})$ (rank). \texttt{False}: find rational points on $C$. Default: \texttt{False}. \\
        FINITE\_FIELD & \texttt{None} for over-$\mathbb{Q}$ search; set to prime $p$ for finite-field DLP mode. Default: \texttt{None}. \\
        SECRET\_KEY & DLP mode: secret exponent $d$ (for benchmarking). \\
        BLOCK\_WIEDEMANN & \texttt{True}: always use Block Wiedemann for the final DLP solve. Default: \texttt{False}. \\
        PRIME\_POOL & Pool of primes for modular residue computation; defaults to \texttt{primes(590)}, auto-filtered per fibration. \\
        NUM\_PRIME\_SUBSETS & Number of independent prime subsets; $\geq 250$ for stability. Default: 500. \\
        USE\_MINIMAL\_MODEL & \texttt{True}: Tate-minimized Weierstrass model (slower, more correct). \texttt{False}: generic fiber model. Default: \texttt{True}. \\
        SYMBOLIC\_SEARCH & \texttt{True}: solve over $\mathbb{Q}(m)$ (usually slower; rarely finds anything extra). Default: \texttt{False}. \\
        HEIGHT\_BOUND & Canonical height cutoff for rational reconstruction candidates. Auto-tuned by default. \\
        DEBUG & Enable verbose diagnostic output. Default: \texttt{False}. \\
        \bottomrule
    \end{tabularx}
\end{table}

\section{Representative Example Outputs}

\subsection{Rank-4 Curve — Full Point Set}

\begin{itemize}[label={\textbullet}]
    \item Curve: $y^2 = x^6 + 8x^5 + 10x^4 - 10x^3 - 11x^2 + 2x + 1$ (LMFDB ``prestige curve'')
    \item Seed point: $(-1, 1)$
    \item \texttt{MUMFORD\_SEARCH} mode: certified rank-4 Mumford basis in 33.5 seconds.
    \item 7 distinct rational $x$-coordinates recovered: $\{0, 1, -3/2, 2, 1/3, -2, -1\}$.
    \item 4 independent Mumford divisors certified by 4 primes, mod-$p$ rank $= 4$.
\end{itemize}

\subsection{Rank-1 Curve — Picard Number Determination}

\begin{itemize}[label={\textbullet}]
    \item Picard number $\rho = 3$ determined exactly (lower = upper $= 3$ at $\ell = 10009$).
    \item Shioda--Tate: $\text{rank}\, E(\mathbb{Q}(m)) = 1$, $\sum (m_v - 1) = 0$.
    \item 2-Selmer bound: $1 \leq \text{rank} \leq 3$.
    \item Total run time: 125 seconds.
\end{itemize}

\subsection{Finite-Field DLP}

\begin{itemize}[label={\textbullet}]
    \item Curve: $y^2 = x^5 + 3x^3 + 2x^2 + 5x + 4$ over $\mathbb{F}_{33554467}$
    \item Secret key $d = 800$, Jacobian order $\approx 10^{12}$, largest prime $\ell \approx 9.7 \times 10^8$.
    \item 18,017 smooth Mumford divisors collected (factor base of 4,504 unique $x$-coordinates).
    \item Relation matrix: $18017 \times 4504$, full rank confirmed. Ready for Block Wiedemann.
    \item Total time: 112 seconds.
\end{itemize}

\subsection{Large-Conductor Genus 2 Curves}

The system has been tested on curves parsed directly from LMFDB database entries (discriminant up to $\sim 10^7$, conductors from hundreds to millions), with \texttt{TERMINATE\_WHEN\_6} targets of 1 to 11. Conductors up to $\sim 10^7$ are handled without special tuning.

\section{Code Architecture}

\begin{table}[h]
    \centering
    \renewcommand{\arraystretch}{1.2}
    \begin{tabularx}{\textwidth}{>{\ttfamily}p{0.38\textwidth} X}
        \toprule
        \textbf{Module} & \textbf{Description} \\
        \midrule
        search7\_genus2.sage & Main entry point: orchestrates fibration tower, search, Mumford basis, analyses. \\
        search\_common.py & Global configuration and shared utilities (curve setup, prime subgroup, CRT helpers); contains LMFDB test curves. \\
        tower.sage & Fibration tower construction: geometry scoring, polynomial interpolation, tower iteration. \\
        search\_main.py & Rational point search: LLL vectors, modular pipeline, CRT, Hensel, rational reconstruction. \\
        mumford\_basis.py & Jacobian basis construction: Mumford CRT reconstruction, torsion filter, mod-$p$ certification. \\
        mumford\_parallel.py & Parallel residue computation for the Mumford system across the prime pool. \\
        mumford\_reconstruction.py & Algebraic pre-filtering and CRT combination of Mumford coordinate residues. \\
        mumford\_core.py & Core Mumford divisor arithmetic and group law. \\
        mumford\_solver.py & Mumford constraint solving over finite fields. \\
        mumford\_equations.py & Mumford condition polynomials. \\
        mumford\_height.py & Height computations for Mumford divisors. \\
        mumford\_verification.py & Verification routines for reconstructed divisors. \\
        index\_calculus.py & Finite-field DLP: factor base generation, relation walk, matrix preparation. \\
        sparse\_linalg\_modp.py & Block Wiedemann sparse linear algebra solver over $\mathbb{Z}/\ell\mathbb{Z}$. \\
        arakelov.py & Arakelov height pairings (finite places) for N\'eron--Tate Gram matrix. \\
        homology.py & Period lattice and archimedean height integrals. \\
        periods.py & Period computation for the Jacobian. \\
        integration.py & Numerical integration utilities. \\
        analytic\_pairings.py & Analytic/archimedean pairing computations. \\
        pairings.py & Height pairing dispatch. \\
        picard.py & Picard number bounds, Lefschetz trace, Shioda--Tate analysis. \\
        selmer.py / selmer\_genus2.py & 2-Selmer rank computation for genus 2 curves. \\
        tate.py & Tate algorithm for Kodaira fiber classification and minimal model reduction. \\
        bounds.py & Splitting field degree estimation, height bound auto-tuning, prime pool management. \\
        ll\_utilities.py & LLL/BKZ lattice reduction, vector generation, rational reconstruction. \\
        heights.py & Canonical height computations. \\
        torsion.py & Torsion subgroup identification. \\
        automorph.py & N\'eron--Severi lattice automorphisms. \\
        sat.py & Saturation analysis for the Mordell--Weil group. \\
        nslattice.py & N\'eron--Severi lattice operations. \\
        local.py / local\_diagnostics.py & Local arithmetic and diagnostics. \\
        diagnostics2.py & Singular fiber analysis. \\
        consensus.py & Multi-fibration consensus filter. \\
        parallel.py & Parallel execution utilities. \\
        core.py & Core shared arithmetic routines. \\
        dedup.py & Deduplication of candidate points. \\
        stats.py & Phase timing, counter tracking, and run statistics. \\
        walker.py & \textbf{Active development} --- Baby-step Giant-step collision walk for DLP relation generation. Experimental; requires \texttt{libwalk.so}. \\
        collision\_walk.c & C source for the collision walk; compile manually to \texttt{libwalk.so}. \\
        \bottomrule
    \end{tabularx}
\end{table}

\section{Known Limitations and Open Questions}

\begin{itemize}[label={\textbullet}]
    \item The 2-Selmer upper bound can greatly exceed the true rank when $\text{Sha}(J)[2]$ is large, giving a wide gap.
    \item Not every rational point is visible from a given fibration: some points may only appear via secondary fibrations or at very high section multiples.
    \item The Mumford basis gives a lower bound on rank, but proving it is the full rank requires the Selmer upper bound to match (or an unconditional height-bound argument).
    \item Local Mumford solutions in characteristic $p$ do not always lift globally; the CRT machinery harvests only those that happen to glue, which remains a theoretical mystery.
    \item Higher genus support: genus 3 rational point search has been tested successfully via the tower construction, but the \texttt{MUMFORD\_SEARCH} and DLP modes currently only support genus 2. Higher genera ($\geq 4$) are largely untested due to difficulty sourcing suitable test curves.
    \item Multi-point fibrations (2+ seed points) are not fully implemented and may hang.
    \item Very high-degree fibrations (deg $> 15$) may require increased \texttt{MAX\_MODULUS}.
    \item The posterior completeness estimator is acknowledged to be unreliable and should not be used to conclude that the search has terminated.
\end{itemize}

\section{Dependencies}

\begin{itemize}[label={\textbullet}]
    \item SageMath $\geq 10.x$ --- core algebraic geometry, Jacobian arithmetic, LLL, CRT.
    \item Python 3.10+ with \texttt{numpy}, \texttt{scipy}, \texttt{tqdm}, \texttt{colorama}, \texttt{cysignals}, \texttt{cypari2}.
    \item PARI/GP (via \texttt{cypari2}) --- polynomial arithmetic and number field computations.
    \item Multiprocessing / \texttt{concurrent.futures} --- parallel residue computation.
    \item Optional: custom C extension (\texttt{collision\_walk.c}), which must be compiled manually to \texttt{libwalk.so} (the binary is not included in the repository). Without it, DLP mode falls back to slower pure-Python walks or raises a \texttt{FileNotFoundError} on import of \texttt{walker.py}.
\end{itemize}

\section{Building This Document}

This document is built with \texttt{latexmk}. A \texttt{Makefile} is provided with the following targets:

\begin{itemize}[label={\textbullet}]
    \item \texttt{make} or \texttt{make all} --- compile \texttt{docs.tex} to \texttt{docs.pdf} using \texttt{pdflatex}.
    \item \texttt{make watch} --- continuous preview mode: recompiles automatically on every save.
    \item \texttt{make clean} --- remove auxiliary files (\texttt{.aux}, \texttt{.log}, etc.), keeping the PDF.
    \item \texttt{make distclean} --- full clean, removing the PDF as well.
\end{itemize}

\noindent\textbf{Prerequisites:} \texttt{latexmk} and a standard \TeX\ distribution (e.g.\ TeX Live or MiKTeX) with \texttt{pdflatex}. All packages used (\texttt{geometry}, \texttt{amsmath}, \texttt{booktabs}, \texttt{tabularx}, \texttt{hyperref}, \texttt{enumitem}, \texttt{titlesec}) are included in a standard full TeX Live installation.

\end{document}
